\origpage{92}{ROBERTSON AND SEYMOUR}
\noindent $E(\Gamma) \ne \emptyset$, \textit{and let $\mathscr{T}$ be a respectful pretangle in $\Gamma$. Let $\Gamma'$ be a radial drawing of $\Gamma$, and let $\mathrm{ins}$ be the slope in $\Gamma'$ derived from $\mathscr{T}$. Then $\mathscr{T}$ is a tangle if and only if $\mathrm{ins}$ is even.} \fix{\textit{The ``if'' part assumes that $\mathscr{T}$ has order $\neq 2$ or that $\Gamma$ has no loop.}}\fixnote{fix:6.5}{As printed, the ``if'' part of (6.5) fails for order 2. Let $\Gamma$ be the drawing in the sphere with one vertex $v$ and one loop $e$. Its radial drawing $\Gamma'$ joins $v$ to the vertex of $\Gamma'$ in each of the two discs bounded by $e$. So $\Gamma'$ has no circuit, and its only region has perimeter of length 4. Hence for $\theta = 2$ there is no circuit or region of length $< 2\theta$, and the slope of order 2 is even. But the induced pretangle of order 2 in $\Gamma$ contains $(\Gamma \setminus e, K_e)$, whose first part is the vertex $v$ alone, and this violates the third axiom of a tangle. With the added hypothesis, (6.5) holds.}

\textit{Proof.} Let $\mathscr{T}$ have order $\theta \ge 1$. Suppose that $\mathscr{T}$ is a tangle and let $r \in R(\Gamma')$ with perimeter of length $< 2\theta$. Let $e \in E(\Gamma)$ with $e \subseteq r$; then $V(\Gamma) \cap \bar{r}$ is the set of ends of $e$, and the perimeter of $r$ has length at least twice the number of ends of $e$, and so $e$ has $< \theta$ ends. Let $K_e$ be the hypergraph formed by $e$ and its ends; then $(\Gamma \setminus e, K_e) \notin \mathscr{T}$ by the third tangle axiom, and so $(K_e, \Gamma \setminus e) \in \mathscr{T}$ by the first. By (6.3) there is a circuit $C$ of $\Gamma'$ with
\begin{equation*}
V(C) \cap V(\Gamma) \subseteq V(K_e \cap (\Gamma \setminus e)) = V(\Gamma) \cap \bar{r}
\end{equation*}
such that $e \subseteq \mathrm{ins}(C)$, that is, $r \subseteq \mathrm{ins}(C)$. Thus $\mathrm{ins}$ is even.

For the converse, suppose that $\mathscr{T}$ is not a tangle. By [4, Theorem (2.7)] there is an edge $e \in E(\Gamma)$ with $< \theta$ ends such that $(\Gamma \setminus e, K_e) \in \mathscr{T}$, with notation as before. Let $r$ be the region of $\Gamma'$ with $e \subseteq r$. For any circuit $C$ of $\Gamma'$ of length $< 2\theta$, $(\Gamma \cap \Delta, \Gamma \cap \overline{\Sigma - \Delta}) \in \mathscr{T}$, where $\Delta = \mathrm{ins}(C)$, and so $(\Gamma \setminus e) \cup (\Gamma \cap \Delta) \ne \Gamma$, by the second tangle axiom; that is, $e \notin \Delta$, or equivalently, $r \nsubseteq \mathrm{ins}(C)$. Hence $\mathrm{ins}$ is not even. $\square$

Another consequence of (6.1) is that, given a drawing in a connected surface, there is a 1-1 correspondence between its respectful tangles and the respectful tangles in its dual. We see this as follows. Let $\Gamma'$ be a 2-cell drawing in a connected surface. One might ask, when is $\Gamma'$ a radial drawing of some 2-cell drawing $\Gamma$ with $E(\Gamma) \ne \emptyset$? First, $\Gamma'$ must be bipartite, so let us assume that $(A, B)$ is a bipartition of $\Gamma'$. Then it is easy to see that $\Gamma'$ is a radial drawing of a 2-cell drawing $\Gamma$ with $V(\Gamma) = A$ and $E(\Gamma) \ne \emptyset$ if and only if every region $r$ of $\Gamma'$ has perimeter of length 4. But if so, then there is another 2-cell drawing $\Gamma^*$, say, with $V(\Gamma^*) = B$, such that $\Gamma'$ is a radial drawing of $\Gamma^*$. We call $\Gamma^*$ the \textit{dual} of $\Gamma$ (it is unique in the usual sense of surface duals). Let $\theta \ge 1$ be an integer. Because there is a 1-1 correspondence between respectful tangles in $\Gamma$ of order $\theta$ and even slopes in $\Gamma'$ of order $\theta$, because of (6.1) and (6.5), and also between respectful tangles in $\Gamma^*$ of order $\theta$ and even slopes in $\Gamma'$ of order $\theta$, for the same reason, we deduce that there is a 1-1 correspondence between the respectful tangles of order $\theta$ in $\Gamma$ and in $\Gamma^*$. In summary:

\stmt
(6.6) \textit{Let $\Gamma$, $\Gamma^*$ be dual 2-cell drawings in a connected surface, with a common radial drawing $\Gamma'$, and let $\theta \ge 1$ be an integer\fix{, such that if $\theta = 2$ then neither $\Gamma$ nor $\Gamma^*$ has a loop}. There is a bijection $\mathscr{T} \leftrightarrow \mathscr{T}^*$ between respectful tangles $\mathscr{T}$ of order $\theta$ in $\Gamma$, and respectful tangles $\mathscr{T}^*$ of order $\theta$ in $\Gamma^*$, such that for each such $\mathscr{T}$, the slopes in $\Gamma'$ derived}
